\section{Introduction}
\label{sec:introduction}

Training language models on their own output across successive generations leads to progressive distribution degradation, a phenomenon variously termed model collapse~\cite{shumailov2024curse} or Model Autophagy Disorder~\cite{alemohammad2024mad}.
Tail modes vanish, lexical diversity narrows, and calibration deteriorates as each generation's output becomes the next generation's training data.
As large language models become ubiquitous content generators, the fraction of machine-generated text on the web grows steadily, making iterative training collapse---where models are unwittingly trained on the output of earlier model generations---a practical concern for the long-term stability of language model ecosystems.

A natural follow-up question is whether cheap-to-compute statistics---\emph{canary metrics}---can provide early warning of impending collapse, analogous to early-warning signals for critical transitions in ecological and climate systems~\cite{scheffer2009early,scheffer2012anticipating}.
Several recent studies have pursued this direction: \citet{du2026geometric} informally observe that expected calibration error rises before accuracy degrades; SIGMA~\cite{sigma2026spectral} proposes spectral indicators of distributional shift; and DCES~\cite{dces2026counter} uses center-edge sampling dynamics as a collapse signal.
Yet the statistical methodology underlying these lead-lag claims---whether a canary metric genuinely \emph{precedes} downstream degradation, or merely correlates with it under shared non-stationary trends---has received little scrutiny.

We identify several methodological deficiencies that pervade existing lead-lag analyses of iterative training collapse:
(1)~no differencing tests to rule out spurious regression arising from shared non-stationary trends---a well-known hazard when correlating co-trending time series;
(2)~no multi-method cross-validation---conclusions rest on a single statistical method whose implicit assumptions about stationarity and directionality are rarely examined;
(3)~no multiple comparison correction (e.g., Benjamini--Hochberg FDR control), increasing the risk of false discoveries when testing many metric pairs simultaneously; and
(4)~no regime-dependent analysis, despite the possibility that the same statistical test yields opposite conclusions at different contamination levels.
The closest predecessor is \citet{du2026geometric}, who observe that calibration error rises before accuracy degrades in iteratively trained language models.
Their analysis, however, is limited to a single metric pair, spans only 1--3 generations, and reports informal visual comparisons without statistical tests---leaving open the question of whether the observed lead is genuine or an artifact of shared trends.
The early-warning framework of \citet{scheffer2009early,scheffer2012anticipating}---widely adopted in ecology for anticipating tipping points---provides the theoretical motivation for canary-based prediction, but has not been systematically applied to iterative training collapse with the methodological rigor that short, noisy time series demand~\cite{boettiger2012limits}.

We address these gaps through a systematic case study of canary prediction of diversity collapse---the progressive loss of lexical diversity in model-generated text across training generations.
Using Pythia-410M across five contamination levels ($\alpha \in \{0.00, 0.25, 0.50, 0.75, 1.00\}$) and up to four random seeds, we expose three methodological pitfalls that undermine naive lead-lag analysis:
\begin{enumerate}
    \item \textbf{Method sensitivity.} Applied to the same data (16 candidate canary$\to$downstream pairs, including candidates later excluded; \S\ref{sec:case_study}), five lead-lag detection methods yield FDR significance rates ranging from 100\% (naive cross-correlation) to 25\% (Toda--Yamamoto Granger causality), with inter-method agreement of Cohen's $\kappa{=}0.185$ (poor agreement).
    \item \textbf{Spurious regression.} Canary$\to$perplexity associations appear near-universal in undifferenced analysis but vanish entirely after first-differencing ($F{\approx}0$, $p{=}0.99$), revealing them as artifacts of co-trending rather than genuine temporal precedence.
    \item \textbf{Regime dependence.} Permutation Granger tests find 24/24 canary$\to$diversity pairs FDR-significant at full synthetic replacement ($\alpha{=}1.0$) but 0/8 at $\alpha{\leq}0.50$, demonstrating that detectability depends sharply on the contamination level.
\end{enumerate}
To move beyond single-method analysis, we propose a multi-method consensus framework with empirically calibrated thresholds, requiring agreement across multiple independent statistical methods and random seeds before classifying a lead-lag relationship as robust.

The framework's utility is demonstrated by its calibration on built-in experimental controls.
Naive cross-correlation identifies 100\% of canary--collapse metric pairs as FDR-significant; the consensus pipeline (requiring ${\geq}4/5$ methods and ${\geq}3/4$ seeds) retains only those pairs independently validated by the first-differencing diagnostic.
On reference extremes, the framework achieves a false positive rate of 0\% at the $\alpha{=}0.0$ null control (where individual methods reach up to 100\% FPR) and a false negative rate of 0\% against six ground-truth pairs at $\alpha{=}1.0$.
Applying the framework, we identify token entropy as a one-generation-ahead canary for diversity collapse at $\alpha{=}1.0$, validated by the full consensus pipeline, with cross-correlation $|r|{>}0.99$ and directional asymmetry confirmed by 0/24 reverse-direction FDR-significant pairs; preliminary single-seed evidence at $\alpha{=}0.75$ awaits multi-seed confirmation.
We note that Granger precedence establishes temporal predictive improvement, not true causality~\cite{shojaie2022granger}.

Our contributions are:
\begin{enumerate}
    \item \textbf{Systematic exposure of three methodological pitfalls} in lead-lag analysis of iterative training collapse.
    We quantify each pitfall: method choice alone swings significance from 100\% to 25\% with a spurious-signal rate of 75\%; first-differencing eliminates canary$\to$perplexity associations entirely; and the same Granger test shifts from 24/24 FDR-significant pairs to 0/8 across contamination levels.
    These results demonstrate that single-method, single-regime analyses---the current norm in this literature---are unreliable.

    \item \textbf{A multi-method consensus framework} with empirically calibrated thresholds and four-tier practical guidance ranging from exploratory ($K{\geq}1/5$) to conservative ($K{=}5/5$).
    The framework achieves FPR${=}$0\% and FNR${=}$0\% on reference extremes, providing robustness against miscalibration of any individual method.

    \item \textbf{Application of the framework to identify token entropy} as a one-generation-ahead canary for diversity collapse at $\alpha{=}1.0$ (with preliminary single-seed evidence at $\alpha{=}0.75$), scoped precisely to diversity metrics (distinct-$n$) rather than model collapse broadly, with a tiered practical guide mapping contamination regimes to recommended analysis strategies.
\end{enumerate}

The remainder of this paper is organized as follows.
\S\ref{sec:related_work} surveys related work on model collapse, early-warning signals, and lead-lag methodology.
\S\ref{sec:background} establishes the formal setup, notation, and the five detection methods used throughout.
\S\ref{sec:method} distills the three methodological pitfalls into a unified diagnostic framework, presents the multi-method consensus pipeline, and reports its empirical calibration.
\S\ref{sec:experiments} provides full experimental results: method sensitivity analysis (\S\ref{sec:method_sensitivity}), regime-dependent detectability (\S\ref{sec:regime}), a case study of canary prediction of diversity collapse (\S\ref{sec:case_study}), and consensus framework validation (\S\ref{sec:consensus_validation}).
\S\ref{sec:discussion} discusses limitations---including the short time-series constraint, single-model scope, and the distinction between Granger precedence and causality---and outlines directions for future work.
\S\ref{sec:conclusion} concludes.